%------------------------------------------------------------------------------%

\usepackage{apresentacao}
\usepackage{mathconfig}
\usepackage{tabto}

\makeatletter
\graphicspath  {{./figs/}}
\def\input@path{{./figs/}}
\makeatother

%------------------------------------------------------------------------------%
\title {Funções de Várias Variáveis}
\author{\large Luis Alberto D'Afonseca}
\date  {Cálculo de Funções de Várias Variáveis I}

%------------------------------------------------------------------------------%
\begin{document}

%------------------------------------------------------------------------------%
\begin{frame}
    \titlepage
\end{frame}

%------------------------------------------------------------------------------%
\section{Funções}

%------------------------------------------------------------------------------%
\begin{frame}{Funções}
  \[
    f\colon A \to B
  \]
  \pause
  \vspace{-\baselineskip}
  \begin{itemize}[<+->]
    \item $A$ \emph{domínio}
    \item $B$ \emph{contradomínio}
    \item $f$ \emph{relação} que associa elementos do domínio a elementos do contradomínio
    \item Cada elemento do domínio deve estar associado a \emph{um único} elemento do contradomínio
    \item A \emph{imagem} é um subconjunto de $B$ composto dos \emph{elementos que são imagem} de algum elemento de $A$
  \end{itemize}
\end{frame}

%------------------------------------------------------------------------------%
\begin{frame}{Domínio}
  \begin{itemize}[<+->]
    \item Precisa ser dado a priori
    \item Geralmente é determinado pela aplicação
    \item Se nada for dito, assumimos o domínio ``natural'', \\
          ie, todos os valores onde podemos efetuar as contas
  \end{itemize}
\end{frame}

%------------------------------------------------------------------------------%
\begin{frame}{Imagem}
  \begin{itemize}[<+->]
    \item Subconjunto do contradomínio
    \item O contradomínio é definido a priori
    \item Pode ser difícil identificar a imagem de uma função
    \item Dada uma função \(f\colon D \to C \) \\
          $y$ pertence a imagem de $f$,
          se existe $x\in D$ tal que
          \[y=f(x)\]
  \end{itemize}
\end{frame}

%------------------------------------------------------------------------------%
\begin{frame}{Nomenclatura -- Curiosidade}
  Função real
  \[
    f\colon D \subset \R \to \R
  \]
  \pause
  Função de várias variáveis
  \[
    f\colon D \subset \R^n \to \R
  \]
  \pause
  Função vetorial ou paramétrica
  \[
    f\colon D \subset \R \to \R^n
  \]
  \pause
  Campo vetorial
  \[
    f\colon D \subset \R^n \to \R^m
  \]
\end{frame}

%------------------------------------------------------------------------------%
\section{Funções de Várias Variáveis}

%------------------------------------------------------------------------------%
\begin{frame}{Funções de Várias Variáveis}
  \[
    f\colon D \subset \R^n \to \R
  \]
  \pause
  \[
    f(x, y) = x + y
  \]
  \pause
  \[
    f(\vec{r}) = \sqrt{\sum_{k=1}^n r_i^2\,}
  \]
\end{frame}


%------------------------------------------------------------------------------%
\framedgraphic{Funções de Várias Variáveis}{definicao_funcao}{}

%------------------------------------------------------------------------------%
\begin{frame}{Domínio e Imagem}

  Determinamos o \emph{domínio} eliminando pontos que levam a divisões por zero ou
  valores complexos.

  \pause
  \vspace{\baselineskip}

  \textit{``A \emph{imagem} consiste on conjunto de valores de saída para a variável dependente''} (sic.)

\end{frame}

%------------------------------------------------------------------------------%
\begin{frame}{Exemplo 1}






















\end{frame}

%------------------------------------------------------------------------------%
\begin{frame}{Intervalos}
  Subconjuntos contínuos de $\R$
  \vspace{\baselineskip}

  \pause
  Intervalo fechado
  \[
    [0, 1] = \left\{ x\in\R\st 0 \leq x \leq 1 \right\}
  \]
  \pause
  Intervalo aberto
  \[
    (-2, 2) = \left\{ x\in\R\st -2 < x < 2 \right\}
  \]
  \pause
  Aberto a esquerda e fechado a direita
  \[
    (1, 5] = \left\{ x\in\R\st 1 < x \leq 5 \right\}
  \]
\end{frame}

%------------------------------------------------------------------------------%
\framedgraphic{Regiões no Plano}{definicao_regioes}{}
\framedgraphic{Regiões no Plano}{pontos_regioes}{}

%------------------------------------------------------------------------------%
\begin{frame}{Regiões Limitadas}
  Uma região do plano é limitada se está dentro de um disco de raio fixo

  \pause
  \vspace{\baselineskip}

  Caso contrário ela é ilimitada
\end{frame}

%------------------------------------------------------------------------------%
\begin{frame}{Exemplo 2}
  Descreva o domínio da função \(f(x, y) = \sqrt{y-x^2}\)

  \pause
  \vspace{\baselineskip}
  \begin{enumerate}[<+->]
    \item Qual a condição que os pontos do domínio deve satisfazer?
    \item Que região é essa? Represente-a graficamente.
    \item Qual o interior da região?
    \item Qual a fronteira da região?
    \item A região é aberta ou fechada?
    \item A região é limitada?
  \end{enumerate}
\end{frame}

%------------------------------------------------------------------------------%
\framedgraphic{Exemplo 2}{dominio_funcao}{}

%------------------------------------------------------------------------------%
\section{Lista Mínima}

%------------------------------------------------------------------------------%
\begin{frame}{Lista Mínima}
  Cálculo Vol. 2 do Thomas 12\textsuperscript{a} ed. -- Seção 14.1
  \begin{enumerate}
    \item Estudar o texto da seção
    \item Resolver os exercícios: 1-2, 5-9
  \end{enumerate}
  \vfill
  Atenção: A prova é baseada no livro, não nas apresentações
\end{frame}

%------------------------------------------------------------------------------%
\begin{frame}
  \tableofcontents
\end{frame}

%------------------------------------------------------------------------------%
\end{document}
%------------------------------------------------------------------------------%
